% §2 Preliminaries

We establish the notation, collect the key definitions used throughout,
and review the standard Majorana decomposition that underpins both
parameterized maps.

\subsection{Fermionic operators and canonical anticommutation relations}

Consider $n$ fermionic modes with creation operators $\ad{j}$ and
annihilation operators $\an{j}$ for $j \in \{0, \ldots, n{-}1\}$.
These satisfy the \emph{canonical anticommutation relations} (CAR):
\begin{align}
  \acomm{\an{i}}{\ad{j}} &= \delta_{ij}, \label{eq:car1} \\
  \acomm{\an{i}}{\an{j}} &= 0, \label{eq:car2} \\
  \acomm{\ad{i}}{\ad{j}} &= 0. \label{eq:car3}
\end{align}
These relations ensure the Pauli exclusion principle: no two fermions
can occupy the same mode.

\subsection{Pauli algebra}

A qubit register of $n$ qubits admits the Pauli group $\mathcal{P}_n$
consisting of $n$-fold tensor products of single-qubit Pauli operators
$\{I, X, Y, Z\}$ with phases in $\{+1, -1, +i, -i\}$.  The
single-qubit multiplication table is:
\begin{equation}
  XY = iZ, \quad YZ = iX, \quad ZX = iY,
  \label{eq:pauli-products}
\end{equation}
with $\sigma^2 = I$ for all $\sigma \in \{X, Y, Z\}$ and
$\sigma_a \sigma_b = -\sigma_b \sigma_a$ for $a \neq b$.
Multi-qubit Pauli strings multiply componentwise:
\begin{equation}
  (\sigma_{a_0} \otimes \cdots \otimes \sigma_{a_{n-1}})
  (\sigma_{b_0} \otimes \cdots \otimes \sigma_{b_{n-1}})
  = \phi \cdot
  (\sigma_{a_0}\sigma_{b_0}) \otimes \cdots \otimes
  (\sigma_{a_{n-1}}\sigma_{b_{n-1}}),
  \label{eq:pauli-mult}
\end{equation}
where $\phi = \prod_{k=0}^{n-1} \phi_k$ is the accumulated phase from
each single-qubit multiplication.  The \emph{Pauli weight} of a string
is the number of non-identity factors.

\subsection{Core definitions}
\label{sec:core-defs}

We collect the central definitions used throughout the paper.  Further
definitions (index-monotonic trees, star trees) are introduced in
\cref{sec:star-tree} where they are needed for the main theorem.

\begin{definition}[Phase group]
\label{def:phase-group}
The \emph{phase group} $\Zi = \{+1, -1, +i, -i\}$ is the cyclic group
of order~4 under multiplication.  We write elements as $\{P_1, M_1,
P_i, M_i\}$.  Every single-qubit Pauli product yields a result Pauli
and a phase element in $\Zi$.
\end{definition}

\begin{definition}[Majorana operators]
\label{def:majorana}
For each mode $j$, the Hermitian \emph{Majorana operators} are:
\begin{align}
  c_j &= \ad{j} + \an{j}, \label{eq:majorana-c} \\
  d_j &= i(\ad{j} - \an{j}). \label{eq:majorana-d}
\end{align}
These satisfy $\{c_i, c_j\} = 2\delta_{ij}$,
$\{d_i, d_j\} = 2\delta_{ij}$, and $\{c_i, d_j\} = 0$.
\end{definition}

\noindent
Inverting~\eqref{eq:majorana-c}--\eqref{eq:majorana-d}:
\begin{equation}
  \ad{j} = \tfrac{1}{2}(c_j - id_j), \qquad
  \an{j} = \tfrac{1}{2}(c_j + id_j).
  \label{eq:ladder-from-majorana}
\end{equation}

\begin{definition}[Fermion-to-qubit encoding]
\label{def:encoding}
A \emph{fermion-to-qubit encoding} for $n$ modes on $n$ qubits is a
pair of maps $c_j \mapsto S_j^{(c)}$ and $d_j \mapsto S_j^{(d)}$, where
each $S_j^{(c)}$ and $S_j^{(d)}$ is a Pauli string (an element of
$\mathcal{P}_n$), such that the induced ladder operators
\begin{equation}
  \ad{j} = \tfrac{1}{2}(S_j^{(c)} - i S_j^{(d)}), \qquad
  \an{j} = \tfrac{1}{2}(S_j^{(c)} + i S_j^{(d)})
  \label{eq:encoded-ladder}
\end{equation}
satisfy the CAR~\eqref{eq:car1}--\eqref{eq:car3}.
\end{definition}

\begin{definition}[Encoding scheme]
\label{def:encoding-scheme}
An \emph{encoding scheme} is a triple of set-valued functions
$\scheme = (\Uset{\cdot}, \Pset{\cdot}, \Occ{\cdot})$ from mode
indices $\{0, \ldots, n{-}1\}$ to subsets of qubit indices
$\{0, \ldots, n{-}1\}$, from which Majorana operators can be
constructed via a universal algorithm (\cref{sec:constructions}).
\end{definition}

\begin{definition}[Combining algebra]
\label{def:combining-algebra}
A \emph{combining algebra} is a function that specifies how to
reorder a ladder operator past another during normal ordering.  The
two instances are:
\begin{itemize}
  \item \emph{Fermionic} (CAR): swapping $\an{j}$ past $\ad{k}$
    with $j = k$ produces two terms with a sign flip
    ($\an{j}\ad{j} = \delta_{jj} - \ad{j}\an{j}$).
  \item \emph{Bosonic} (CCR): swapping $b_j$ past $b_k^\dagger$
    with $j = k$ produces two terms with no sign change
    ($b_j b_j^\dagger = b_j^\dagger b_j + 1$).
\end{itemize}
\end{definition}

\noindent
The central observation is that every known encoding can be specified
by describing how to construct $S_j^{(c)}$ and $S_j^{(d)}$, and two
distinct parameterizations accomplish this.

\subsection{The electronic Hamiltonian}

The second-quantized electronic Hamiltonian is:
\begin{equation}
  H = \sum_{pq} h_{pq}\, \ad{p}\an{q}
    + \frac{1}{2} \sum_{pqrs} g_{pqrs}\, \ad{p}\ad{q}\an{s}\an{r},
  \label{eq:hamiltonian}
\end{equation}
where $h_{pq}$ are one-electron integrals and $g_{pqrs}$ are
two-electron integrals from quantum chemistry
computations~\cite{szabo1996, helgaker2000}.  An encoding maps each
$\ad{j}$ and $\an{j}$ to a two-term Pauli expression
via~\eqref{eq:encoded-ladder}, and the products are expanded using the
Pauli multiplication rule~\eqref{eq:pauli-mult}.  The result is a
Hamiltonian expressed entirely in terms of Pauli strings:
\begin{equation}
  H = \sum_{\alpha} h_\alpha\, \sigma_\alpha,
  \label{eq:pauli-hamiltonian}
\end{equation}
where $\sigma_\alpha \in \mathcal{P}_n$ and $h_\alpha \in \mathbb{C}$.
